% GENERATED FILE. Edit canonical lessons, then run npm run generate:books.
% canonical-source-hash: fnv1a64:732b4b502f87298e
% canonical-lessons: TE-C371-mantrikudu, TE-C371-maya, TE-C371-deyyam, TE-C371-varam, TE-C371-sapam

\chapter{Magic and Ghosts}
\label{ch:te-magic-and-ghosts}
\hlchaptermodality{371}

\begin{chapteropening}
\textbf{By the end of this chapter:} \emph{I can say a sorcerer, magic, a ghost, a boon and a curse.}

Complete the last lesson of chapter 371: I can say a sorcerer, magic, a ghost, a boon and a curse.
\end{chapteropening}

\section[\texorpdfstring{\te{మాంత్రికుడు} (māntrikuḍu)}{మాంత్రికుడు (māntrikuḍu)}]{\te{మాంత్రికుడు} (māntrikuḍu) — a sorcerer; a magician}

\label{lesson:TE-C371-mantrikudu}



\begin{quote}
Before the new one: say the Telugu for to extol, then the Telugu for to chastise.
\end{quote}

\subsection*{You'll want to know: \te{మాంత్రికుడు}}
\textbf{\te{మాంత్రికుడు}} — \emph{māntrikuḍu} — \textquotedblleft{}a sorcerer; a magician\textquotedblright{}.

\begin{quote}
\te{మాంత్రికుడు} \te{రాకుమారిని} \te{బంధించాడు}.
\end{quote}

\begin{quote}
\emph{māntrikuḍu rākumārini bandhiñcāḍu.}
\end{quote}

\begin{quote}
The sorcerer imprisoned the princess.
\end{quote}

Plural: \te{మాంత్రికులు}, \emph{māntrikulu}.

From \te{మంత్రం}, a spell: one who works spells.

\begin{cousinweb}
\textbf{\te{మాంత్రికుడు}} (\emph{māntrikuḍu}) is borrowed from Sanskrit \emph{māntrika}, \textquotedblleft{}a reciter of spells, a sorcerer\textquotedblright{} (CDIAL 10050), from \emph{mantra}, \textquotedblleft{}a spell\textquotedblright{}.
\end{cousinweb}

\begin{grammarlens}[title={What you've built}]
One more word for telling what happened.
\end{grammarlens}

\subsection*{Your turn}
\begin{itemize}
\raggedright
  \item \emph{Say it:} \emph{māntrikuḍu}
  \item \emph{Say it:} \emph{māntrikuḍu}, once more
  \item \emph{Say it:} \emph{māntrikuḍu rākumārini bandhiñcāḍu.} — The sorcerer imprisoned the princess.
  \item \emph{Recall:} say the Telugu for to extol, then the Telugu for to chastise, then say \emph{māntrikuḍu} again
\end{itemize}

\begin{tcolorbox}[breakable,colback=teal!4,colframe=teal!35!black,title={Before you move on}]
What does \te{మాంత్రికుడు} mean? (\textquotedblleft{}A sorcerer; a magician\textquotedblright{}.) Say it once more.
\end{tcolorbox}
\section[\texorpdfstring{\te{మాయ} (māya)}{మాయ (māya)}]{\te{మాయ} (māya) — magic; an illusion}

\label{lesson:TE-C371-maya}



\begin{quote}
Before the new one: say the Telugu for to chastise, then the Telugu for a sorcerer.
\end{quote}

\subsection*{You'll want to know: \te{మాయ}}
\textbf{\te{మాయ}} — \emph{māya} — \textquotedblleft{}magic; an illusion\textquotedblright{}.

\begin{quote}
\te{మాంత్రికుడు} \te{మాయతో} \te{ఒక} \te{కోట} \te{సృష్టించాడు}.
\end{quote}

\begin{quote}
\emph{māntrikuḍu māyatō oka kōṭa sṛṣṭiñcāḍu.}
\end{quote}

\begin{quote}
The sorcerer created a castle by magic.
\end{quote}

\te{మాయ} \te{చేయు}, to work magic; also a trick.

\begin{cousinweb}
\textbf{\te{మాయ}} (\emph{māya}) is borrowed from Sanskrit \emph{māyā}, \textquotedblleft{}supernatural power, skill\textquotedblright{}, and \textquotedblleft{}illusion\textquotedblright{} (CDIAL 10060).
\end{cousinweb}

\begin{grammarlens}[title={What you've built}]
One more word for telling what happened.
\end{grammarlens}

\subsection*{Your turn}
\begin{itemize}
\raggedright
  \item \emph{Say it:} \emph{māya}
  \item \emph{Say it:} \emph{māya}, once more
  \item \emph{Say it:} \emph{māntrikuḍu māyatō oka kōṭa sṛṣṭiñcāḍu.} — The sorcerer created a castle by magic.
  \item \emph{Recall:} say the Telugu for to chastise, then the Telugu for a sorcerer, then say \emph{māya} again
\end{itemize}

\begin{tcolorbox}[breakable,colback=teal!4,colframe=teal!35!black,title={Before you move on}]
What does \te{మాయ} mean? (\textquotedblleft{}Magic; an illusion\textquotedblright{}.) Say it once more.
\end{tcolorbox}
\section[\texorpdfstring{\te{దెయ్యం} (deyyaṁ)}{దెయ్యం (deyyaṁ)}]{\te{దెయ్యం} (deyyaṁ) — a ghost}

\label{lesson:TE-C371-deyyam}



\begin{quote}
Before the new one: say the Telugu for a sorcerer, then the Telugu for magic.
\end{quote}

\subsection*{You'll want to know: \te{దెయ్యం}}
\textbf{\te{దెయ్యం}} — \emph{deyyaṁ} — \textquotedblleft{}a ghost\textquotedblright{}.

\begin{quote}
\te{ఆ} \te{పాత} \te{ఇంట్లో} \te{దెయ్యం} \te{ఉందని} \te{పిల్లలు} \te{అనుకుంటారు}.
\end{quote}

\begin{quote}
\emph{ā pāta iṇṭlō deyyaṁ undani pillalu anukuṇṭāru.}
\end{quote}

\begin{quote}
The children think there is a ghost in that old house.
\end{quote}

Plural: \te{దెయ్యాలు}, \emph{deyyālu}.

\begin{cousinweb}
\textbf{\te{దెయ్యం}} (\emph{deyyaṁ}) is usually traced to Sanskrit \emph{daitya}, \textquotedblleft{}a demon\textquotedblright{} (Monier-Williams).
\end{cousinweb}

\begin{grammarlens}[title={What you've built}]
One more word for telling what happened.
\end{grammarlens}

\subsection*{Your turn}
\begin{itemize}
\raggedright
  \item \emph{Say it:} \emph{deyyaṁ}
  \item \emph{Say it:} \emph{deyyaṁ}, once more
  \item \emph{Say it:} \emph{ā pāta iṇṭlō deyyaṁ undani pillalu anukuṇṭāru.} — The children think there is a ghost in that old house.
  \item \emph{Recall:} say the Telugu for a sorcerer, then the Telugu for magic, then say \emph{deyyaṁ} again
\end{itemize}

\begin{tcolorbox}[breakable,colback=teal!4,colframe=teal!35!black,title={Before you move on}]
What does \te{దెయ్యం} mean? (\textquotedblleft{}A ghost\textquotedblright{}.) Say it once more.
\end{tcolorbox}
\section[\texorpdfstring{\te{వరం} (varaṁ)}{వరం (varaṁ)}]{\te{వరం} (varaṁ) — a boon; a blessing}

\label{lesson:TE-C371-varam}



\begin{quote}
Before the new one: say the Telugu for magic, then the Telugu for a ghost.
\end{quote}

\subsection*{You'll want to know: \te{వరం}}
\textbf{\te{వరం}} — \emph{varaṁ} — \textquotedblleft{}a boon; a blessing\textquotedblright{}.

\begin{quote}
\te{మాంత్రికుడు} \te{రాజుకు} \te{ఒక} \te{వరం} \te{ఇచ్చాడు}.
\end{quote}

\begin{quote}
\emph{māntrikuḍu rājuku oka varaṁ iccāḍu.}
\end{quote}

\begin{quote}
The sorcerer granted the king a boon.
\end{quote}

Plural: \te{వరాలు}, \emph{varālu}.

Also in daily speech: \te{అది} \te{నాకు} \te{వరం}, that is a blessing for me.

\begin{cousinweb}
\textbf{\te{వరం}} (\emph{varaṁ}) is borrowed from Sanskrit \emph{vara}, \textquotedblleft{}choice, a boon\textquotedblright{} (CDIAL 11308).
\end{cousinweb}

\begin{grammarlens}[title={What you've built}]
One more word for telling what happened.
\end{grammarlens}

\subsection*{Your turn}
\begin{itemize}
\raggedright
  \item \emph{Say it:} \emph{varaṁ}
  \item \emph{Say it:} \emph{varaṁ}, once more
  \item \emph{Say it:} \emph{māntrikuḍu rājuku oka varaṁ iccāḍu.} — The sorcerer granted the king a boon.
  \item \emph{Recall:} say the Telugu for magic, then the Telugu for a ghost, then say \emph{varaṁ} again
\end{itemize}

\begin{tcolorbox}[breakable,colback=teal!4,colframe=teal!35!black,title={Before you move on}]
What does \te{వరం} mean? (\textquotedblleft{}A boon; a blessing\textquotedblright{}.) Say it once more.
\end{tcolorbox}
\section[\texorpdfstring{\te{శాపం} (śāpaṁ)}{శాపం (śāpaṁ)}]{\te{శాపం} (śāpaṁ) — a curse}

\label{lesson:TE-C371-sapam}



\begin{quote}
Before the new one: say the Telugu for a ghost, then the Telugu for a boon.
\end{quote}

\subsection*{You'll want to know: \te{శాపం}}
\textbf{\te{శాపం}} — \emph{śāpaṁ} — \textquotedblleft{}a curse\textquotedblright{}.

\begin{quote}
\te{ఆ} \te{శాపం} \te{వల్ల} \te{రాకుమారుడు} \te{రాయిగా} \te{మారాడు}.
\end{quote}

\begin{quote}
\emph{ā śāpaṁ valla rākumāruḍu rāyigā mārāḍu.}
\end{quote}

\begin{quote}
Because of that curse, the prince turned to stone.
\end{quote}

Plural: \te{శాపాలు}, \emph{śāpālu}.

\te{శాపం} \te{పెట్టు}, to lay a curse.

\begin{cousinweb}
\textbf{\te{శాపం}} (\emph{śāpaṁ}) is borrowed from Sanskrit \emph{śāpa}, \textquotedblleft{}a curse, abuse\textquotedblright{} (CDIAL 12393).
\end{cousinweb}

\begin{grammarlens}[title={What you've built}]
One more word for telling what happened.
\end{grammarlens}

\subsection*{Your turn}
\begin{itemize}
\raggedright
  \item \emph{Say it:} \emph{śāpaṁ}
  \item \emph{Say it:} \emph{śāpaṁ}, once more
  \item \emph{Say it:} \emph{ā śāpaṁ valla rākumāruḍu rāyigā mārāḍu.} — Because of that curse, the prince turned to stone.
  \item \emph{Recall:} say the Telugu for a ghost, then the Telugu for a boon, then say \emph{śāpaṁ} again
\end{itemize}

\begin{tcolorbox}[breakable,colback=teal!4,colframe=teal!35!black,title={Before you move on}]
What does \te{శాపం} mean? (\textquotedblleft{}A curse\textquotedblright{}.) Say it once more.
\end{tcolorbox}
